\chapter{Introduction}

\section{Contexte du projet}

MSI Consulting est une entreprise fictive de taille moyenne (50–100 employés) répartie en six départements :
\begin{itemize}
    \item \textbf{Direction} : 5 postes de cadres supérieurs
    \item \textbf{Commercial} : 3 postes de vendeurs et managers commerciaux
    \item \textbf{Informatique} : 3 techniciens IT et administrateurs
    \item \textbf{Ressources Humaines} : 2 responsables RH
    \item \textbf{Comptabilité} : 2 responsables financiers
    \item \textbf{Serveurs} : infrastructure centralisée (DCs, NAS)
\end{itemize}

L'infrastructure est hébergée sur \textbf{VMware Workstation Pro} et comporte :
\begin{itemize}
    \item 1 pare-feu \textbf{pfSense} (routeur et filtrage)
    \item 2 contrôleurs de domaine \textbf{Windows Server 2022}
    \item 1 serveur de fichiers centralisé
    \item 6 réseaux virtuels isolés (VLANs)
    \item 1 workstation de test
\end{itemize}

Le domaine est \texttt{corpmsi.local}.

\section{Objectifs pédagogiques}

Ce projet couvre l'intégralité des compétences requises pour administrer une infrastructure Active Directory en entreprise :

\begin{enumerate}
    \item \textbf{Active Directory Domain Services}
    \begin{itemize}
        \item Déploiement d'une forêt et domaine
        \item Création de Domain Controllers en réplication
        \item Gestion des OUs et utilisateurs
    \end{itemize}

    \item \textbf{Services réseau critiques}
    \begin{itemize}
        \item DNS : zones directives et inverses
        \item DHCP : étendues, failover, DHCP Relay
        \item NTP : synchronisation temps sur PDC
    \end{itemize}

    \item \textbf{Politique de sécurité}
    \begin{itemize}
        \item Politiques de groupe (GPOs) par département
        \item Politiques de mot de passe strictes (PSO)
        \item BitLocker et chiffrement de disque
        \item Contrôle d'accès aux fichiers (NTFS)
    \end{itemize}

    \item \textbf{Infrastructure de fichiers}
    \begin{itemize}
        \item Partages SMB par département
        \item DFS Namespaces pour transparence réseau
        \item Énumération basée sur l'accès (ABE)
    \end{itemize}

    \item \textbf{Pare-feu et segmentation réseau}
    \begin{itemize}
        \item pfSense : configuration multi-VLAN
        \item Règles de pare-feu par interface
        \item Isolation inter-VLAN sauf pour IT
        \item Filtrage applicatif
    \end{itemize}

    \item \textbf{Proxy et filtrage URL}
    \begin{itemize}
        \item Squid en mode transparent
        \item SquidGuard avec blacklists (UT1)
        \item Inspection SSL MITM
        \item Politiques de filtrage par département
    \end{itemize}

    \item \textbf{Troubleshooting}
    \begin{itemize}
        \item Diagnostic avec \texttt{dcdiag}
        \item Analyse DNS avec \texttt{nslookup}
        \item Vérification de réplication AD
        \item Résolution de problèmes courants
    \end{itemize}
\end{enumerate}

\section{Compétences développées}

À la fin de ce projet, vous maîtriserez :

\begin{description}
    \item[Administration Active Directory] Configuration d'OUs, GPOs, PSO, replication inter-DC
    \item[Gestion réseau] VLANs, segmentation, pare-feu stateful, DHCP Relay
    \item[Sécurité] Chiffrement BitLocker, politiques de mot de passe, contrôle d'accès
    \item[Diagnostic] Identification et résolution rapide de problèmes AD/DNS
    \item[Infrastructure réelle] Compréhension de productions réelles en entreprise
    \item[PowerShell] Scripts d'automatisation pour déploiement et maintenance
\end{description}

\section{Ce que j'ai appris}

Durant la mise en place de ce projet, plusieurs découvertes majeures ont émergé :

\begin{enumerate}
    \item \textbf{L'importance du PDC}
    \begin{itemize}
        \item Le PDC Emulator doit être impérativement le serveur NTP fiable
        \item Sans synchronisation temps correcte, Kerberos échoue silencieusement
        \item Les timezone mal configurées causent des comportements anormaux subtils
    \end{itemize}

    \item \textbf{DNS et AD sont inséparables}
    \begin{itemize}
        \item AD repose entièrement sur DNS (SRV records, zone directives/inverses)
        \item Un problème DNS minimal (ex : records GC manquants) peut paralyser partiellement le domaine
        \item Les records sont créés automatiquement par Netlogon, mais peuvent nécessiter correction manuelle
    \end{itemize}

    \item \textbf{Le secure channel entre DCs}
    \begin{itemize}
        \item Le canal sécurisé Kerberos entre DC et domaine peut se casser
        \item La commande \texttt{nltest /dsregdns} est essentielle pour diagnostiquer
        \item \texttt{Nltest.exe /sc\_change\_pwd:domain} répare le canal en forçant le renouvellement du mot de passe machine
    \end{itemize}

    \item \textbf{VLANs et segmentation réalistes}
    \begin{itemize}
        \item La segmentation VLAN ne doit pas être prise à la légère
        \item Bloquer les ports AD inadéquatement paralyse les clients
        \item L'IT doit avoir un accès transversal complet, les autres VLANs doivent être isolés
    \end{itemize}

    \item \textbf{Les ports RPC dynamiques sont un cauchemar}
    \begin{itemize}
        \item Par défaut : ports 49152–65535 (16 383 ports !)
        \item Restreindre à 60000–60099 permet des règles pare-feu gérables
        \item Indispensable pour GPO, DFS, replication AD
    \end{itemize}

    \item \textbf{Le proxy SSL MITM complexe}
    \begin{itemize}
        \item Dépendre d'une CA personnalisée crée des dépendances GPO complexes
        \item La whitelist SSL (No Bump) est critique pour les services de paiement et mises à jour
        \item Sans GPO de déploiement du CA, tous les navigateurs affichent des avertissements
    \end{itemize}

    \item \textbf{BitLocker et TPM}
    \begin{itemize}
        \item BitLocker obligatoire nécessite TPM sur tous les postes (ou PIN supplémentaire)
        \item La sauvegarde des clés dans AD via GPO est essentielle pour la récupération
        \item Sans cela, un poste perdu = données perdues
    \end{itemize}
\end{enumerate}

\section{Structure du document}

Ce document est organisé comme suit :

\begin{description}
    \item[Chapitre 2 — Architecture] Plan d'adressage, VLANs, OUs, assignation réseau VMware
    \item[Chapitre 3 — pfSense] Installation, configuration, règles firewall, NAT
    \item[Chapitre 4 — GPOs] Toutes les politiques de groupe détaillées par département
    \item[Chapitre 5 — DFS et fichiers] Serveur de fichiers, partages, DFS Namespaces
    \item[Chapitre 6 — Proxy] Squid, SquidGuard, SSL MITM, ACLs
    \item[Chapitre 7 — NTP] Configuration du serveur NTP sur PDC
    \item[Chapitre 8 — Troubleshooting] Cas d'étude DNS complet avec diagnostique et résolution
    \item[Chapitre 9 — Améliorations] Futures extensions du projet
\end{description}

\section{Prérequis pour le lecteur}

Pour comprendre ce document, vous devez maîtriser :

\begin{itemize}
    \item Les bases d'Active Directory (OUs, GPOs, utilisateurs)
    \item La configuration réseau IPv4 (subnets, masques)
    \item L'administration Windows Server via GUI et PowerShell
    \item Les concepts de pare-feu et filtering réseau
\end{itemize}

Les commandes PowerShell et les fichiers de configuration sont fournis dans le dépôt \texttt{/scripts} et \texttt{/configs}.
